\documentclass[11pt]{article}

\usepackage[margin=1in]{geometry}
\usepackage{amsmath,amssymb,mathtools}
\usepackage{bm}
\usepackage{enumitem}
\usepackage{hyperref}
\usepackage{xcolor}

\hypersetup{
  colorlinks=true,
  linkcolor=blue,
  citecolor=blue,
  urlcolor=blue
}

\newcommand{\R}{\mathbb{R}}
\newcommand{\C}{\mathbb{C}}
\newcommand{\A}{\mathcal{A}}
\newcommand{\D}{\mathcal{D}}
\newcommand{\G}{\mathcal{G}}
\newcommand{\V}{\mathcal{V}}
\newcommand{\norm}[1]{\left\lVert #1 \right\rVert}
\newcommand{\ip}[2]{\left\langle #1,#2 \right\rangle}

\title{Born-Coordinate ATFPS Preconditioning for mgRTE}
\author{mgRTE technical note}
\date{\today}

\begin{document}
\maketitle

\begin{abstract}
This note gives the conclusion and detailed rationale for using a
Born-coordinate learned preconditioner in the current ATFPS discretization of
the steady radiative transfer equation (RTE).  The conclusion is that the
safest and most useful transfer from the Neural Preconditioned Born Series
(NPBS) Helmholtz code is not to replace ATFPS by a pseudo-spectral Born solver.
Instead, the current ATFPS block inverse \texttt{funcDinv} should be used as
the RTE analogue of the Helmholtz Green/reference inverse.  The neural network
should be trained and evaluated in this preconditioned residual coordinate, and
FGMRES should remain right-preconditioned with the true ATFPS residual as the
stopping criterion.
\end{abstract}

\section{Conclusion}

Let \(A_h\) denote the current matrix-free ATFPS operator implemented by
\texttt{func\_inflow\_Dirichlet}.  Let \(D_h\) denote its local block-diagonal
part implemented by \texttt{funcD}, and let
\[
  G_D := D_h^{-1}
\]
be the local inverse implemented by \texttt{funcDinv}.  The proposed
Born-coordinate ATFPS preconditioner is
\[
  q = G_D r,\qquad z_\theta = P_\theta(q;\sigma_T,\sigma_a,\log\varepsilon),
\]
followed by the ATFPS consistency correction
\[
  M_\theta(r) = \operatorname{Corr}(z_\theta,r_{\rm unselected}).
\]
This \(M_\theta\) is then used as a right preconditioner inside FGMRES:
\[
  A_h M_\theta y = b,\qquad x = M_\theta y.
\]

The key conclusion is:
\begin{quote}
\emph{For the current mgRTE repository, Born-coordinate ATFPS is the right
near-term idea.  It preserves the ATFPS discretization, keeps the existing
boundary/layer handling, and imports the most important NPBS idea: train the
network in the same reference-inverse coordinate in which it will be used by
Krylov acceleration.}
\end{quote}

This is mathematically more defensible than claiming a drop-in Helmholtz-style
Convergent Born Series for RTE.  For RTE, a standalone stationary Born iteration
has no simple scalar damping factor analogous to Helmholtz
\(\gamma=iV/\eta\), and diffusion/interface regimes can be close to
noncontractive.  Using the Born/reference inverse as a right preconditioner in
FGMRES is the safer formulation because convergence is monitored using the true
residual \(\norm{b-A_h x}\).

\section{The ATFPS Discrete Operator}

After discrete ordinates in angle, a scaled steady RTE can be written
schematically as
\begin{equation}
  c_m \partial_x \psi_m
  + s_m \partial_y \psi_m
  + \frac{\sigma_T}{\varepsilon}\psi_m
  - \left(\frac{\sigma_T}{\varepsilon}
          - \varepsilon\sigma_a\right)
    \sum_n \kappa_{mn}\omega_n \psi_n
  = \varepsilon q_m.
  \label{eq:dom-rte}
\end{equation}
The ATFPS method used in mgRTE constructs local exponential basis functions in
each cell.  In the code this is visible in the directional blocks
\begin{equation}
  B_x =
  \operatorname{diag}(1/c_m)
  \left[
    \left(1-\varepsilon^2\frac{\sigma_a}{\sigma_T}\right)K W
    - I
  \right],
  \qquad
  B_y =
  \operatorname{diag}(1/s_m)
  \left[
    \left(1-\varepsilon^2\frac{\sigma_a}{\sigma_T}\right)K W
    - I
  \right],
  \label{eq:directional-blocks}
\end{equation}
where \(K=(\kappa_{mn})\) and \(W=\operatorname{diag}(\omega_n)\).  The
functions \texttt{VDXY}, \texttt{FSM}, \texttt{Inflow2alpha},
\texttt{MLRO\_Dirichlet}, and \texttt{MBTO\_Dirichlet} assemble the local
basis, boundary-to-coefficient maps, and left/right/bottom/top interface maps.

The unknown used by the current ATFPS solver is not the angular flux
\(\psi\in\R^{4M I J}\).  It is an inflow/interface basis vector
\[
  u \in \R^{8M I J}.
\]
The global residual operator is
\begin{equation}
  A_h u = r,
  \label{eq:atfps-operator}
\end{equation}
implemented by \texttt{func\_inflow\_Dirichlet}.  The selected ATFPS modes are
determined by the tensor \texttt{VecSize}; selected and unselected components
are separated by \texttt{restrict\_basis} and \texttt{squeeze\_basis}.

It is useful to view the global ATFPS operator as
\begin{equation}
  A_h = D_h - E_h,
  \label{eq:block-splitting}
\end{equation}
where \(D_h\) contains the local cell/interface blocks and \(E_h\) contains
neighbor transfer terms.  The code does not explicitly store this matrix, but
the split is exactly represented by
\[
  D_h u = \texttt{funcD}(u),
  \qquad
  D_h^{-1}r = \texttt{funcDinv}(r).
\]
This is the important object for Born-coordinate training.

\section{Why \texttt{funcDinv} Is the RTE Born Coordinate}

The NPBS Helmholtz implementation uses a reference Green operator \(G_\eta\) to
map residuals into Born coordinates:
\[
  q = G_\eta r.
\]
The network is then trained and evaluated in that same coordinate.  The
analogous ATFPS choice is
\[
  G_D = D_h^{-1}.
\]
This is not a formal free-space Green function, but it is a better
discrete-reference inverse for the current repository for three reasons.

\begin{enumerate}[leftmargin=2em]
  \item \textbf{It matches the implemented unknown.}
  The network and FGMRES currently operate on ATFPS inflow/interface variables
  \(u\in\R^{8M I J}\), not on angular flux values.  A Fourier or continuous
  Green function would naturally live on angular flux variables
  \(\psi\in\R^{4M I J}\), so using it directly would require a separate
  discretization backend.

  \item \textbf{It captures local stiff scales.}
  In diffusion regimes, \(\varepsilon\ll 1\), the local exponential modes and
  filtered basis selection are the numerically delicate part of the method.
  \(D_h^{-1}\) uses the same local ATFPS blocks and therefore respects this
  scaling.

  \item \textbf{It preserves boundary and compression logic.}
  ATFPS boundary conditions and selected/unselected mode consistency are built
  into the current operator.  \texttt{funcDinv} and \texttt{correction} work in
  that same geometry, whereas a naive FFT Green operator would impose periodic
  structure and ignore inflow boundary treatment.
\end{enumerate}

Thus the Born coordinate for ATFPS should be
\begin{equation}
  q = G_D r = D_h^{-1} r.
  \label{eq:born-coordinate}
\end{equation}
The learned preconditioner should receive \(q\), not \(r\), whenever it is
trained as a Born-coordinate model.

\section{Training Objectives}

The original raw-residual objective trains a model \(P_\theta\) by minimizing
\begin{equation}
  \mathcal{L}_{\rm raw}(\theta)
  =
  \norm{
    S(A_h P_\theta(r)-r)
  }_2^2,
  \label{eq:raw-loss}
\end{equation}
where \(S\) extracts the selected ATFPS modes.  This is a reasonable baseline,
but it has a coordinate mismatch: the model must learn both the local stiff
inverse \(D_h^{-1}\) and the remaining global correction.

The Born-coordinate objective first maps the residual into the local inverse
coordinate:
\begin{equation}
  q = D_h^{-1}r.
\end{equation}
Then the network output is interpreted as a candidate solution correction
\begin{equation}
  z_\theta = P_\theta(q;\sigma_T,\sigma_a,\log\varepsilon).
\end{equation}
The primary Born-coordinate loss is
\begin{equation}
  \mathcal{L}_{D}(\theta)
  =
  \frac{
    \norm{
      S\left(D_h^{-1}A_h z_\theta - q\right)
    }_2^2
  }{
    \norm{S q}_2^2 + \delta
  },
  \label{eq:born-loss}
\end{equation}
where \(\delta>0\) is a small numerical safeguard.

The mixed loss used for the recommended experiment is
\begin{equation}
  \mathcal{L}_{\rm mixed}(\theta)
  =
  \mathcal{L}_{D}(\theta)
  +
  \lambda
  \frac{
    \norm{S(A_h z_\theta-r)}_2^2
  }{
    \norm{S r}_2^2 + \delta
  }.
  \label{eq:mixed-loss}
\end{equation}
The second term prevents the network from only looking good in the
\(D_h^{-1}\)-metric while still producing a poor true ATFPS residual.  In the
current implementation \(\lambda=0.1\) is used as a conservative default.

\subsection{Direct and Residual Forms}

There are two useful network interpretations.

\paragraph{Direct form.}
\begin{equation}
  z_\theta = P_\theta(q).
\end{equation}
This keeps the old model contract: the network directly returns a correction.

\paragraph{Residual-over-reference form.}
\begin{equation}
  z_\theta = q + P_\theta(q).
  \label{eq:residual-form}
\end{equation}
This is safer if one wants the learned model to start from block Jacobi and
learn only the nonlocal correction.  If \(P_\theta\approx0\), then the method
falls back to \(D_h^{-1}\).  The code supports both forms through
\texttt{preconditioner\_form=direct} and \texttt{preconditioner\_form=residual}.

\section{FGMRES Evaluation}

The learned preconditioner is nonlinear and may depend on coefficient features,
so it should be used as a flexible right preconditioner.  FGMRES stores
\[
  z_j = M(v_j)
\]
for each Arnoldi vector \(v_j\), and reconstructs
\[
  x_k = x_0 + \sum_{j=0}^{k-1} y_j z_j.
\]
The stopping criterion must remain the true residual
\begin{equation}
  \frac{\norm{b-A_h x_k}_2}{\norm{b}_2}.
  \label{eq:true-residual}
\end{equation}

For Born-coordinate ATFPS, the preconditioner used by FGMRES is
\begin{align}
  q_j &= D_h^{-1}v_j,\\
  \widetilde z_j &= P_\theta(q_j),\\
  z_j &= \operatorname{Corr}(\widetilde z_j,(v_j)_{\rm unselected}).
  \label{eq:fgmres-born-preconditioner}
\end{align}
The correction step is important because selected-only training can otherwise
hide errors in unselected ATFPS modes.  The code implements this pathway in
\texttt{rte\_preconditioning.py} and exposes it in inference as
\texttt{mgnet\_born}.  The original raw-residual MgNet mode remains available
as \texttt{mgnet\_raw}.

The evaluation should compare at least
\[
  I,\qquad D_h^{-1},\qquad M_{\theta,{\rm raw}},\qquad M_{\theta,{\rm born}}.
\]
Both the true residual \eqref{eq:true-residual} and the reference-coordinate
metric
\begin{equation}
  \frac{\norm{D_h^{-1}(b-A_h x_k)}_2}{\norm{D_h^{-1}b}_2}
  \label{eq:d-metric}
\end{equation}
should be reported.

\section{Relation to a Convergent Born Series}

The block splitting \eqref{eq:block-splitting} yields the fixed-point equation
\begin{equation}
  u = D_h^{-1}(r + E_h u).
\end{equation}
The corresponding stationary iteration is
\begin{equation}
  u^{(n+1)} = D_h^{-1}r + D_h^{-1}E_h u^{(n)}.
  \label{eq:block-born-iteration}
\end{equation}
It converges if
\begin{equation}
  \rho(D_h^{-1}E_h) < 1
  \quad\text{or, more conservatively,}\quad
  \norm{D_h^{-1}E_h}<1.
  \label{eq:contraction}
\end{equation}
This condition may hold in benign transport/high-absorption regimes, but it is
not guaranteed in diffusion or interface regimes.  In the diffusion limit,
scattering ratios approach one and transport/source iteration analogues become
slow or nearly noncontractive.  Therefore the correct claim is not that
Born-coordinate ATFPS gives a guaranteed standalone Convergent Born Series.
The correct claim is that \(D_h^{-1}\) is a strong reference inverse and a
natural coordinate for a learned right preconditioner.

This distinction matters.  Helmholtz CBS uses a complex shifted reference
operator and a scalar damping factor designed around the Helmholtz resolvent.
RTE is dissipative, anisotropic in direction, and boundary driven.  Its
contraction behavior is governed by absorption, scattering, boundary leakage,
and diffusion scaling.  FGMRES is therefore the robust outer iteration.

\section{Relation to Pseudo-Spectral RTE}

A pseudo-spectral RTE reference operator can be useful for a separate research
backend.  For periodic or padded smooth problems, a constant-coefficient DOM
reference has Fourier symbol
\begin{equation}
  \widehat A_0(\xi)
  =
  i\,\operatorname{diag}(\Omega_m\cdot \xi)
  +
  C_0,
  \label{eq:pseudospectral-symbol}
\end{equation}
where \(C_0\) is a constant angular collision matrix.  The reference inverse is
\begin{equation}
  G_0 f =
  \mathcal{F}^{-1}
  \left[
    \widehat A_0(\xi)^{-1}\widehat f(\xi)
  \right].
  \label{eq:fft-green}
\end{equation}
This is implemented experimentally in \texttt{RTE\_pseudospectral.py}.

However, this operator lives naturally on angular flux states
\(\psi\in\C^{4M I J}\), not on the ATFPS inflow basis
\(u\in\R^{8M I J}\).  It also imposes periodic/padded structure unless extra
boundary treatment is added.  Therefore it should not replace ATFPS for the
current diffusion/interface benchmarks.  It is best viewed as a separate
smooth-problem reference or an FFT preconditioner experiment.

\section{Implementation Summary}

The current repository implements the following concrete design.

\begin{itemize}[leftmargin=2em]
  \item \texttt{rte\_preconditioning.py}: shared ATFPS preconditioner helpers,
  including raw and Born-coordinate MgNet application.

  \item \texttt{train.py}: training accepts
  \texttt{loss\_type=raw}, \texttt{loss\_type=born}, and
  \texttt{loss\_type=born\_mixed}.

  \item \texttt{inference.py}: inference can compare
  \texttt{block\_jacobi}, \texttt{mgnet\_raw}, and \texttt{mgnet\_born}
  under right-preconditioned FGMRES.

  \item \texttt{RTE\_mgmodel.py}: \texttt{MG\_precond} now supports custom
  input/output channels, needed for future \(4M\)-channel pseudo-spectral
  states while preserving the default \(8M\)-channel ATFPS behavior.

  \item \texttt{RTE\_pseudospectral.py}: an experimental periodic
  pseudo-spectral DOM operator with matrix-valued FFT Green inverse.
\end{itemize}

\section{Synthetic Comparison Test}

A small synthetic block-diagonal ATFPS test was added to verify the coordinate
idea.  In that test, the operator is constructed so that \(D_h^{-1}\) is the
exact inverse.  An identity ``model'' is then compared in raw and Born
coordinates.  The observed result was
\[
  \mathcal{L}_{\rm raw}=1.149097,
  \qquad
  \mathcal{L}_{\rm born\_mixed}=5.64\times 10^{-19}.
\]
This test does not prove performance on real heterogeneous RTE instances, but
it verifies the core mechanism: when the reference inverse is the correct
coordinate, the Born-coordinate objective removes the local scaling error that
the raw objective forces the network to learn.

The pseudo-spectral reference test also verifies the separate FFT idea on a
constant-coefficient periodic problem:
\[
  \text{identity residual history}
  =
  (1.0,\;9.53\times 10^{-1},\;7.45\times 10^{-1},\;6.89\times 10^{-1}),
\]
while the FFT Green preconditioner gives
\[
  \text{FFT Green residual history}
  =
  (1.0,\;2.89\times 10^{-7}).
\]
Again, this validates the reference inverse on its natural problem class, not
on the full ATFPS boundary-value problem.

\section{Recommended Experiments}

The recommended next experiments are:

\begin{enumerate}[leftmargin=2em]
  \item Generate the existing ATFPS datasets for diffusion and interface
  regimes.

  \item Train a raw baseline with \texttt{configs/default.yaml}.

  \item Train a Born-coordinate model with \texttt{configs/born\_precond.yaml}.

  \item Evaluate
  \[
    D_h^{-1},\quad M_{\theta,{\rm raw}},\quad M_{\theta,{\rm born}}
  \]
  using the same FGMRES tolerance, restart, random seed, and test coefficients.

  \item Report iteration count, final true residual, final \(D_h^{-1}\)-metric
  residual, and wall-clock cost.

  \item Add replay residuals from partial block-Jacobi or FGMRES runs to the
  training set if random residuals do not transfer well to solver residuals.
\end{enumerate}

\section{Risks and Failure Modes}

\begin{itemize}[leftmargin=2em]
  \item \textbf{Train/inference coordinate mismatch.}
  A Born-trained model must receive \(D_h^{-1}r\) at inference time.  Feeding
  raw \(r\) into a Born-trained model invalidates the training objective.

  \item \textbf{Selected-only loss can hide unselected-mode errors.}
  The correction step should remain in both training and FGMRES
  preconditioning.

  \item \textbf{\(D_h^{-1}\) can be ill-conditioned.}
  Diffusion/interface regimes may produce large \(\norm{D_h^{-1}r}\).  Training
  should monitor NaNs, local solve failures, and the ratio
  \(\norm{D_h^{-1}r}/\norm{r}\).

  \item \textbf{Random residuals may be insufficient.}
  A model trained only on random right-hand sides may not see the residual
  distribution generated by FGMRES.  Replay residuals are a likely improvement.

  \item \textbf{Compression threshold dependence.}
  Changing \texttt{tol\_exp} changes \texttt{VecSize} and therefore the
  selected/unselected mode geometry.  Checkpoints should be tied to the
  compression threshold used for training.
\end{itemize}

\section{Final Position}

Born-coordinate ATFPS should be treated as the main learned-preconditioner
direction for the current mgRTE codebase.  It is a conservative refactor:
it preserves the arXiv:2604.23265-style ATFPS discretization, respects inflow
boundary handling and filtered modes, and imports the central NPBS lesson that
the network should be trained in a reference-inverse residual coordinate.

The pseudo-spectral/FFT Green operator is still valuable, but it belongs in a
parallel research backend.  It is suitable for smooth periodic or padded
benchmarks and for studying matrix-valued Green preconditioners.  It should not
be advertised as a replacement for ATFPS in diffusion/interface regimes until
boundary treatment, asymptotic-preserving behavior, and interface robustness
are demonstrated.

\begin{thebibliography}{9}

\bibitem{mgRTE-paper}
arXiv:2604.23265.
The ATFPS/MgNet radiative transfer equation work referenced by this repository.
\url{https://arxiv.org/abs/2604.23265}.

\bibitem{npbs-code}
NeuralPreconditionedBornSeries Helmholtz reference implementation.
Local reference path:
\texttt{/Users/hxchen/NeuralPreconditionedBornSeries/helmholtz}.

\end{thebibliography}

\end{document}
